\documentclass[11pt,a4paper]{article}

\usepackage[margin=0.9in]{geometry}
\usepackage{fontspec}
\usepackage{xeCJK}
\xeCJKsetup{CJKspace=true}
\setmainfont{DejaVu Serif}
\setsansfont{DejaVu Sans}
\setmonofont{DejaVu Sans Mono}
\setCJKmainfont{Noto Sans CJK KR}
\setCJKsansfont{Noto Sans CJK KR}
\setCJKmonofont{Noto Sans Mono CJK KR}
\usepackage{amsmath,amssymb,mathtools}
\usepackage{booktabs}
\usepackage{longtable}
\usepackage{tabularx}
\usepackage{array}
\usepackage{enumitem}
\usepackage{xcolor}
\usepackage{hyperref}
\usepackage{listings}
\usepackage{titlesec}
\usepackage{fancyhdr}
\usepackage{microtype}
\usepackage{ragged2e}

\hypersetup{colorlinks=true, linkcolor=blue!55!black, urlcolor=blue!55!black, citecolor=blue!55!black}
\definecolor{notegray}{RGB}{245,246,248}
\definecolor{darkblue}{RGB}{30,58,138}
\definecolor{darkgreen}{RGB}{22,101,52}
\definecolor{darkred}{RGB}{153,27,27}
\definecolor{softblue}{RGB}{239,246,255}

\lstset{basicstyle=\ttfamily\small, breaklines=true, columns=fullflexible, frame=single, backgroundcolor=\color{notegray}, xleftmargin=0.5em, xrightmargin=0.5em, framerule=0pt}

\setlength{\headheight}{14pt}
\emergencystretch=3em
\pagestyle{fancy}
\fancyhf{}
\lhead{3DGS Wind Response Pipeline}
\rhead{2026-05-12}
\cfoot{\thepage}

\titleformat{\section}{\Large\bfseries\color{darkblue}}{\thesection}{0.8em}{}
\titleformat{\subsection}{\large\bfseries}{\thesubsection}{0.8em}{}
\titleformat{\subsubsection}{\normalsize\bfseries}{\thesubsubsection}{0.8em}{}

\newcommand{\clamp}{\operatorname{clamp}}
\newcommand{\smoothstep}{\operatorname{smoothstep}}
\newcommand{\normalize}{\operatorname{normalize}}
\newcommand{\Var}{\operatorname{Var}}
\newcommand{\GeoMean}{\operatorname{GeoMean}}
\newcommand{\argmin}{\operatorname*{arg\,min}}
\newcommand{\Rraw}{\widetilde{R}_G}
\newcommand{\eps}{\varepsilon}
\newcommand{\todo}[1]{\textcolor{darkred}{#1}}
\newcommand{\phase}[1]{\textcolor{darkgreen}{\textbf{#1}}}
\newcommand{\refbased}{\textcolor{darkblue}{\textbf{[Ref]}}}
\newcommand{\ours}{\textcolor{darkgreen}{\textbf{[Ours]}}}
\newcommand{\heuristic}{\textcolor{darkred}{\textbf{[Heuristic]}}}

\newcolumntype{Y}{>{\RaggedRight\arraybackslash}X}
\newcolumntype{L}[1]{>{\RaggedRight\arraybackslash}p{#1}}

\title{\textbf{Persistent Mesh-Proxy Wind Response Pipeline for 3DGS}\\
\large Module Order, Operations, Effects, and Runtime Flow\\
\normalsize Pipeline record for \textit{Wind-Driven Deformable 3D Gaussian Splatting}}
\author{Prepared as a repo-level pipeline note}
\date{May 12, 2026}

\begin{document}
\maketitle

\begin{center}
\fbox{\parbox{0.92\linewidth}{
\textbf{Pipeline name.} \textbf{Persistent Mesh-Proxy Wind Response Pipeline for 3DGS}.\\[0.35em]
\textbf{One-line summary.} Given a trained 3DGS asset, construct a fixed-topology mesh-proxy scaffold, bind each Gaussian persistently to the proxy, compute wind loads on fixed material-space load samples, solve proxy dynamics, transport mesh deformation back to Gaussian attributes, and render the dynamic 3DGS response.\\[0.35em]
\textbf{Central loop.}
\[
\begin{aligned}
\mathrm{3DGS}&\rightarrow \mathrm{simulation\ proxy}
\rightarrow \mathrm{wind\ dynamics}\\
&\rightarrow \mathrm{Gaussian\ attribute\ transport}
\rightarrow \mathrm{3DGS\ rendering}.
\end{aligned}
\]
}}
\end{center}

\section{Purpose of This Note}

이 문서는 현재 작성된 아이디어 스케치를 기준으로, 입력 3DGS가 들어온 뒤 최종 response, 즉 dynamic rendered frame 또는 video가 출력될 때까지 필요한 모듈을 순서대로 정리한다. 목적은 contribution 문장을 다시 쓰기 전에, 실제 pipeline dependency를 명확히 하는 것이다.

중요한 원칙은 다음과 같다.

\begin{itemize}[leftmargin=1.4em]
    \item 모든 모듈이 runtime에 매 프레임 실행되는 것은 아니다. Asset setup, training/preprocessing, runtime module을 분리해야 한다.
    \item 논문의 중심은 모든 구성요소의 단순 조합이 아니라, 3DGS와 topology-aware wind simulation 사이의 \textbf{persistent representation bridge}이다.
    \item 3DGS는 최종 input과 rendering representation으로 남고, mesh proxy는 내부 simulation scaffold로 쓰인다.
    \item anchor, load sample, Gaussian binding은 모두 고정된 material-space correspondence를 만들기 위한 장치이다. Runtime에는 topology나 binding을 바꾸지 않는다.
\end{itemize}

\section{Pipeline-Level Input and Output}

\subsection{Input}

기본 입력은 trained 3D Gaussian Splatting asset이다.

\begin{equation}
G^0 = \{g_j^0\}_{j=1}^{N_G},\qquad
 g_j^0=(\mu_j^0,R_j^0,s_j^0,\alpha_j^0,c_j^0).
\end{equation}

여기서 각 항은 단순한 이름이 아니라, trained 3DGS가 갖고 있는 canonical rendering state를 의미한다.
일반적인 Inria-style 3DGS PLY에서는 scale, opacity, rotation이 unconstrained parameter로 저장되어 있으므로, loader에서 물리적으로 유효한 값으로 변환해 사용한다.

\begin{center}
\small
\begin{tabularx}{0.98\linewidth}{L{0.12\linewidth}L{0.35\linewidth}Y}
\toprule
항 & 의미 & 값의 계산 또는 획득 방법 \\
\midrule
$\mu_j^0$ &
Canonical Gaussian center. Gaussian $g_j$의 기준 시점 중심 위치이며, proxy extraction, Gaussian-to-proxy binding, covariance transport의 기준 좌표가 된다. &
PLY의 position field, 보통 $(x_j,y_j,z_j)$에서 직접 읽는다. 이 값은 trained 3DGS의 canonical object/world coordinate에 있으며, 이후 deformation은 이 위치를 직접 최적화하지 않고 proxy binding을 통해 갱신한다. \\
\addlinespace
$R_j^0$ &
Canonical Gaussian orientation. Anisotropic ellipsoid의 local axes가 world/canonical coordinate에서 어떤 방향을 향하는지 나타낸다. &
PLY의 $(\hat q_{j0},\hat q_{j1},\hat q_{j2},\hat q_{j3})$ 또는 \texttt{rot\_0..3} quaternion을 normalize한 뒤 rotation matrix로 변환한다:
\[
R_j^0=\operatorname{Rot}\left(\frac{\hat q_j}{\|\hat q_j\|+\eps}\right).
\]
이 값은 covariance 계산과 triangle-local frame transport의 초기 orientation으로 사용된다. \\
\addlinespace
$s_j^0$ &
Canonical anisotropic scale. Gaussian ellipsoid의 세 principal-axis standard deviation 또는 radius 역할을 한다. &
PLY의 \texttt{scale\_0..2}는 보통 log-scale unconstrained parameter $\hat s_j$로 저장된다. 양수 scale은
\[
s_j^0=\exp(\hat s_j)
\]
로 계산한다. 이후 covariance는
\[
\Sigma_j^0=R_j^0\operatorname{diag}((s_j^0)^2)(R_j^0)^\top
\]
로 구성한다. \\
\addlinespace
$\alpha_j^0$ &
Canonical opacity. Rendering alpha contribution, floater filtering, Gaussian reliability, wind-axis transmittance의 기본 opacity evidence로 쓰인다. &
PLY의 \texttt{opacity}는 보통 logit parameter $\hat\alpha_j$로 저장된다. 실제 opacity는
\[
\alpha_j^0=\frac{1}{1+\exp(-\hat\alpha_j)}
\]
로 계산하며, low-opacity Gaussian은 floater likelihood나 reliability score에서 낮은 신뢰도를 갖는다. \\
\addlinespace
$c_j^0$ &
Canonical appearance coefficients. View-dependent color를 표현하는 SH coefficients이며, runtime에서는 SH residual correction의 기준 appearance가 된다. &
PLY의 \texttt{f\_dc\_*}와 선택적 \texttt{f\_rest\_*}에서 읽는다. \texttt{f\_dc}는 SH degree 0 color component이고, higher-order SH가 있으면 $c_j^0$에 함께 stack한다. Debug color는 보통
\[
\mathrm{rgb}_{j}^{\mathrm{dc}}\approx \operatorname{clip}(C_0 f_{\mathrm{dc},j}+0.5,0,1)
\]
로 확인하지만, 최종 rendering은 full SH coefficients를 renderer convention에 맞게 사용한다. \\
\bottomrule
\end{tabularx}
\normalsize
\end{center}

위 항들을 이용해 preprocessing 단계에서 다음 파생량을 만든다.

\begin{equation}
\Sigma_j^0=R_j^0\operatorname{diag}((s_j^0)^2)(R_j^0)^\top,
\qquad
n_j^G=\operatorname{eigvec}_{\min}(\Sigma_j^0).
\end{equation}

$\Sigma_j^0$는 Gaussian covariance transport의 기준이고, $n_j^G$는 가장 얇은 축 방향으로부터 얻은 surface-normal 후보이다.
이 normal은 proxy surface consistency, floater likelihood, local-frame binding의 보조 evidence로 사용된다.

추가 입력은 사용자 또는 데이터 기반 wind signal, constraint, 그리고 rendering camera sequence이다.

\begin{equation}
W(x,t),\qquad \mathcal{C}_{\mathrm{user}},\qquad \Pi_t.
\end{equation}

\begin{center}
\small
\begin{tabularx}{0.98\linewidth}{L{0.16\linewidth}L{0.34\linewidth}Y}
\toprule
항 & 의미 & 값의 계산 또는 획득 방법 \\
\midrule
$W(x,t)$ &
공간과 시간에 따라 변하는 wind velocity field. Load sample 위치에서 effective wind를 계산하는 입력이다. &
MVP에서는 analytic one-way field로 둔다:
\[
\begin{aligned}
W(x,t)=&\,W_{\mathrm{global}}(t)
+\sum_k W_{\mathrm{source},k}(x,t)\\
&+W_{\mathrm{gust}}(x,t).
\end{aligned}
\]
즉 global direction/strength, fan source, sinusoidal gust, turbulence noise parameter로 계산한다. 이후 단계에서는 measured wind 또는 learned latent wind로 대체할 수 있다. \\
\addlinespace
$\mathcal{C}_{\mathrm{user}}$ &
User constraint 또는 editable control field. 어떤 부분이 고정, 부착, 핸들, 높은 stiffness, crease/rib인지 지정한다. &
사용자 paint mask, profile preset, boundary selection, geodesic distance field, manually selected vertices/faces에서 만든다. Solver에서는 inverse mass, attachment compliance, pinned vertex mask, material stiffness field로 변환한다. \\
\addlinespace
$\Pi_t$ &
Rendering camera 또는 evaluation view sequence. Dynamic Gaussian state $G(t)$를 어떤 view에서 render할지 정한다. &
각 frame의 camera intrinsics/extrinsics로 구성된다. 예를 들어 synthetic setup에서는 JSON camera file의 $K$, camera-to-world 또는 world-to-camera matrix에서 읽는다. Training/evaluation에서는 ground-truth view와 같은 camera path를 사용하고, preview에서는 turntable 또는 fixed camera를 사용한다. \\
\bottomrule
\end{tabularx}
\normalsize
\end{center}

\subsection{Output}

최종 출력은 time-varying Gaussian asset과 rendered response이다.

\begin{equation}
G(t)=\{g_j(t)\}_{j=1}^{N_G},\qquad
I_t=\mathcal{R}_{\mathrm{3DGS}}(G(t),\Pi_t).
\end{equation}

영상 출력이면:

\begin{equation}
\mathcal{I}=\{I_1,I_2,\ldots,I_T\}.
\end{equation}

내부 상태로는 deformed proxy mesh가 존재한다.

\begin{equation}
M(t)=(V(t),E,F).
\end{equation}

$M(t)$는 사용자가 직접 보는 output은 아니지만, wind deformation, boundary condition, force propagation, Gaussian transport를 연결하는 핵심 상태이다.

\section{Phase Separation}

전체 pipeline은 세 phase로 나누는 것이 가장 명확하다.

\begin{center}
\begin{tabularx}{0.96\linewidth}{L{0.2\linewidth}L{0.18\linewidth}Y}
\toprule
Phase & 실행 시점 & 포함 모듈 \\
\midrule
Asset setup & asset당 1회 & 3DGS preprocessing, proxy extraction, simulation proxy construction, anchor/load sample setup, Gaussian binding \\
Training / preprocessing & dataset 또는 asset 준비 중 & Blender dense supervision, dense-to-proxy force projection, projection-risk-based load sample refinement, force/appearance model training \\
Runtime response & 매 frame & wind evaluation, exposure estimation, force-space load prediction, load-to-anchor projection, XPBD solver, Gaussian attribute transport, SH correction, rendering \\
\bottomrule
\end{tabularx}
\end{center}

핵심적으로, \textbf{runtime에는 remeshing, Gaussian rebinding, load sample 재배치가 없다}. Runtime에서는 고정된 topology, 고정된 binding, 고정된 material-space samples 위에서 state와 force value만 갱신한다.

\section{Module Order Overview}

\begin{longtable}{L{0.07\linewidth}L{0.25\linewidth}L{0.18\linewidth}L{0.39\linewidth}}
\toprule
순서 & 모듈 & 실행 시점 & 핵심 역할 \\
\midrule
\endfirsthead
\toprule
순서 & 모듈 & 실행 시점 & 핵심 역할 \\
\midrule
\endhead
1 & 3DGS preprocessing & Asset setup & floater 완화, surface-oriented attribute와 Gaussian reliability 후보 계산 \\
2 & Proxy mesh extraction & Asset setup & 3DGS에서 raw topology scaffold 생성 \\
3 & Simulation-oriented proxy construction & Asset setup & render mesh가 아니라 compact physics mesh 구성 \\
4 & Anchor and load sample setup & Asset setup & simulation anchors와 fixed material-space load samples 생성 \\
5 & Gaussian-to-proxy binding & Asset setup & Gaussian을 triangle id, barycentric coordinate, local frame에 persistently binding \\
6 & Blender dense supervision & Training only & dense force, deformation, rendered reference 생성 \\
7 & Dense-to-proxy force projection & Training / preprocessing & dense force를 proxy/load target으로 압축하면서 net force와 wrench 보존 \\
8 & Runtime wind field evaluation & Runtime & 현재 load sample 위치에서 wind 계산 \\
9 & Wind exposure estimation & Runtime & Gaussian/proxy opacity 기반 exposure와 effective wind 계산 \\
10 & Force-space load prediction & Runtime & deformation state가 아니라 aerodynamic load 예측 \\
11 & Load-to-anchor projection & Runtime & load sample force를 simulation vertices로 누적 \\
12 & Proxy dynamics solver & Runtime & XPBD/PBD로 topology-aware deformation 계산 \\
13 & Gaussian attribute transport & Runtime & mesh deformation을 Gaussian mean/covariance/local frame으로 전달 \\
14 & SH / appearance correction & Runtime & analytic SH correction과 residual로 appearance 보정 \\
15 & 3DGS rendering output & Runtime & updated Gaussian set을 standard 3DGS renderer로 출력 \\
\bottomrule
\end{longtable}

\section{Module 1: 3DGS Preprocessing}

\subsection*{Input}

\begin{equation}
G^0=\{g_j^0\}_{j=1}^{N_G},\qquad
g_j^0=(\mu_j^0,R_j^0,s_j^0,\alpha_j^0,c_j^0).
\end{equation}

이 식에서 $G^0$는 wind deformation을 적용하기 전의 전체 trained 3DGS object 또는 asset을 의미한다.
$g_j^0$는 그 object를 구성하는 $j$번째 Gaussian splat 하나이다.
따라서 $G^0$는 하나의 연속 surface mesh가 아니라, $N_G$개의 Gaussian state가 모인 집합으로 본다.

각 $g_j^0$는 독립적인 작은 rendering primitive의 canonical state이다.
$\mu_j^0$는 중심 위치, $R_j^0$는 local ellipsoid axes의 방향, $s_j^0$는 anisotropic scale, $\alpha_j^0$는 opacity, $c_j^0$는 color 또는 SH appearance coefficients를 나타낸다.
위첨자 $0$는 이 값들이 deformed state가 아니라 rest/canonical state에서 읽힌 값임을 뜻한다.
Runtime wind response에서는 이 per-Gaussian state가 proxy binding을 통해 $g_j(t)$로 갱신된다.

\noindent\refbased\;
이 입력 표현은 기존 3D Gaussian Splatting의 표준 state representation에 기반한다 \cite{kerbl2023gaussian}.
논문 작성 시에는 이 식 자체를 novelty로 주장하지 않고, trained 3DGS asset을 읽는 기본 형식으로 둔다.
우리의 관심은 이후 단계에서 이 state를 wind-oriented proxy construction과 persistent binding의 evidence로 어떻게 사용하는지에 있다.

\subsection*{Operation}

먼저 각 Gaussian의 anisotropic shape을 하나의 matrix로 표현하기 위해 Gaussian covariance를 계산한다.
$s_j^0$는 local ellipsoid axes 방향의 scale이고, $R_j^0$는 그 local axes를 canonical/world coordinate로 회전시키는 orientation이다.
따라서 diagonal covariance를 만든 뒤 $R_j^0$로 회전시키면, Gaussian $g_j^0$가 3D 공간에서 실제로 어떤 방향으로 길고 얇게 퍼져 있는지 표현할 수 있다.

\begin{equation}
\Sigma_j^0 = R_j^0\operatorname{diag}((s_j^0)^2)(R_j^0)^\top.
\end{equation}

\noindent\refbased\;
이 covariance 계산은 Gaussian scale과 rotation에서 covariance를 구성하는 표준적인 Gaussian parameter conversion이다.
3DGS는 anisotropic covariance를 갖는 3D Gaussians를 scene primitive로 사용하므로, 이 구성은 기존 3DGS 표현에 기반한다 \cite{kerbl2023gaussian}.
이후 transport의 기준으로 쓰이지만, covariance 구성 자체는 기존 Gaussian 수학 및 3DGS 표현에 속한다.

여기서 scale을 제곱하는 이유는 covariance가 standard deviation 자체가 아니라 variance를 담는 matrix이기 때문이다.
이 $\Sigma_j^0$는 이후 Gaussian covariance transport의 기준이 된다.
즉 proxy mesh가 휘거나 회전할 때 Gaussian center만 옮기는 것이 아니라, Gaussian ellipsoid의 방향과 퍼짐도 surface deformation을 따라가도록 만들기 위한 초기 shape descriptor이다.

다음으로 가장 작은 scale 방향을 Gaussian normal 후보로 사용할 수 있다.
Surface 근처의 anisotropic Gaussian은 보통 surface tangent 방향으로 넓게 퍼지고, surface normal 방향으로 얇게 퍼진다.
따라서 covariance의 가장 작은 eigenvalue에 해당하는 eigenvector는 이 Gaussian이 암시하는 얇은 방향, 즉 surface normal 후보로 해석할 수 있다.

\begin{equation}
n_j^G=\operatorname{eigvec}_{\min}(\Sigma_j^0).
\end{equation}

\noindent\refbased\;+\;\ours\;
Gaussian orientation과 surface alignment를 활용하는 관점은 SuGaR나 GOF 같은 surface-aware Gaussian/mesh extraction 계열 연구와 연결된다 \cite{guedon2024sugar,yu2024gof}.
다만 여기서는 $n_j^G$를 확정 normal로 주장하지 않고, proxy extraction과 Gaussian-to-proxy binding을 안정화하기 위한 보조 evidence로 쓰는 adaptation이다.

다만 $n_j^G$는 확정적인 mesh normal이 아니라 Gaussian에서 얻은 후보 normal이다.
거의 isotropic한 Gaussian, noisy Gaussian, floater Gaussian에서는 이 방향이 불안정할 수 있고, eigenvector의 부호도 양/음이 모호하다.
그래서 이 값은 proxy extraction이나 binding에서 surface consistency를 판단하는 보조 evidence로 사용하고, 이후 proxy mesh normal 또는 local frame과 함께 검증한다.

마지막으로 opacity reliability를 계산한다.
Opacity $\alpha_j^0$가 낮은 Gaussian은 실제 surface를 안정적으로 설명하기보다 floater, background noise, weak contribution일 가능성이 상대적으로 크다.
하지만 opacity의 절대값 분포는 asset마다 다르기 때문에 고정 threshold를 쓰지 않고, asset 내부의 opacity percentile을 이용해 정규화한다.

\begin{equation}
q_j^\alpha=\smoothstep(P_{\mathrm{low}}(\alpha^0),P_{\mathrm{high}}(\alpha^0),\alpha_j^0).
\end{equation}

\noindent\heuristic\;+\;\ours\;
이 항은 표준 3DGS 식이라기보다 asset-wise opacity 분포 차이를 줄이기 위한 robust heuristic이다.
$q_j^\alpha$는 본 논문의 main contribution이 아니라, proxy extraction과 binding에서 unreliable Gaussian의 영향력을 낮추기 위한 supporting design으로 둔다.

$P_{\mathrm{low}}(\alpha^0)$와 $P_{\mathrm{high}}(\alpha^0)$는 전체 asset의 canonical opacity set $\{\alpha_j^0\}_{j=1}^{N_G}$에서 계산한 low/high percentile threshold이다.
이 percentile normalization의 목적은 asset마다 달라지는 opacity scale에 대해 robust한 상대 신뢰도를 만드는 것이다.
따라서 특정 숫자 20/80 자체가 물리 법칙이나 3DGS 이론에서 고정되는 값은 아니며, MVP에서 사용할 수 있는 heuristic default로 본다.
초기 구현에서는
\[
P_{\mathrm{low}}=P_{20},\qquad P_{\mathrm{high}}=P_{80}
\]
로 시작하고, 이후 $P_{10}/P_{90}$, $P_{20}/P_{80}$, $P_{30}/P_{70}$ 같은 ablation으로 proxy extraction과 binding 안정성에 미치는 영향을 확인한다.
$\smoothstep$은 낮은 opacity 영역을 0에 가깝게, 높은 opacity 영역을 1에 가깝게 부드럽게 보낸다.
따라서 $q_j^\alpha$는 Gaussian을 즉시 삭제하기 위한 binary mask가 아니라, 이 Gaussian이 surface evidence로 얼마나 믿을 만한지 나타내는 continuous reliability weight이다.
MVP에서는 이 값을 proxy extraction, binding score, floater likelihood에서 공통 confidence로 사용할 수 있다.

\subsection*{Output}

\begin{equation}
\widetilde{G}^0=\{(g_j^0,q_j^\alpha,n_j^G)\}_{j=1}^{N_G}.
\end{equation}

\vspace{-0.5em}

\noindent\ours\;
$\widetilde{G}^0$로 원본 state와 보조 evidence를 묶어 넘기는 것은 pipeline-specific design이다.
이는 새 rendering representation을 제안하는 것이 아니라, wind-oriented proxy extraction과 persistent Gaussian-to-proxy binding을 안정화하기 위한 preprocessing package이다.

$\widetilde{G}^0$는 원본 3DGS asset을 바꾼 새로운 rendering asset이라기보다, preprocessing을 통해 보조 정보를 붙인 enriched Gaussian set이다.
각 Gaussian은 원래 state $g_j^0$를 그대로 유지하면서, reliability $q_j^\alpha$와 Gaussian normal candidate $n_j^G$를 추가로 갖는다.
이 출력은 다음 단계인 proxy mesh extraction과 Gaussian-to-proxy binding에서 사용된다.

\subsection*{Effect}

이 단계의 효과는 raw 3DGS를 곧바로 mesh extraction이나 binding에 넣지 않고, 각 Gaussian이 얼마나 surface-like한지와 어떤 local orientation을 암시하는지 먼저 평가한다는 점이다.
Low-opacity floater나 noisy Gaussian은 proxy extraction과 binding을 망칠 수 있으므로 $q_j^\alpha$를 통해 영향력을 낮춘다.
반대로 신뢰도 높은 anisotropic Gaussian은 covariance와 normal 후보를 통해 surface orientation의 단서를 제공한다.

주의할 점은 preprocessing이 aggressive pruning 단계가 아니라는 것이다.
특히 isolated Gaussian을 직접 제거하면 leaf tip, antenna-like thin structure, sparse valid detail처럼 실제로 의미 있는 얇은 구조가 사라질 수 있다.
따라서 isolation은 단독 제거 조건으로 쓰지 않고, low opacity, surface distance, user protection score 등과 함께 gated floater likelihood에서만 사용한다.

\section{Module 2: Proxy Mesh Extraction}

\subsection*{Input}

\begin{equation}
\widetilde{G}^0.
\end{equation}

$\widetilde{G}^0$는 Module 1에서 원본 Gaussian state에 reliability와 normal candidate를 붙인 enriched Gaussian set이다.
이 모듈은 그 Gaussian 집합을 바로 simulation에 쓰지 않고, 먼저 connected surface representation으로 바꾸는 단계이다.
입력 3DGS는 point-like Gaussian들의 집합이므로 edge, face, boundary, geodesic distance 같은 topological relation이 명시적으로 존재하지 않는다.

\subsection*{Operation}

3DGS에서 raw surface 또는 mesh proxy를 추출한다.

\begin{equation}
\widetilde{G}^0\rightarrow M_{\mathrm{raw}}=(V_{\mathrm{raw}},E_{\mathrm{raw}},F_{\mathrm{raw}}).
\end{equation}

$M_{\mathrm{raw}}$는 extraction method가 만든 초기 mesh이다.
$V_{\mathrm{raw}}$는 vertex positions, $E_{\mathrm{raw}}$는 edges, $F_{\mathrm{raw}}$는 triangular faces를 의미한다.

본 프로젝트의 기본 선택은 GOF-style extraction을 main raw proxy extraction으로 두는 것이다 \cite{yu2024gof}.
GOF는 trained 3DGS asset에서 opacity field를 통해 surface를 복원하는 방향이므로, 현재 pipeline의 입력을 3DGS 그대로 유지하면서 mesh proxy를 얻기에 가장 자연스럽다.
따라서 real trained 3DGS asset에 대한 기본 경로는 다음처럼 둔다.

\begin{equation}
\widetilde{G}^0 \xrightarrow{\mathcal{E}_{\mathrm{GOF}}} M_{\mathrm{raw}}.
\end{equation}

\noindent\refbased\;+\;\ours\;
GOF-style extraction 자체는 기존 3DGS-to-surface reconstruction baseline으로 사용한다.
우리의 novelty는 GOF mesh extraction 알고리즘을 새로 제안하는 데 있지 않고, 추출된 rendering-oriented raw mesh를 wind simulation을 위한 compact proxy로 바꾸고, 그 위에 persistent Gaussian binding, fixed load samples, anchor-based dynamics, SH correction을 연결하는 데 있다.

비교군으로는 SuGaR-style extraction과 TSDF-style extraction을 둔다.
SuGaR는 Gaussian을 surface-aligned하게 regularize하고 mesh reconstruction과 mesh rendering을 함께 고려하므로 mesh editing이나 surface correspondence 관점에서 강한 비교군이 될 수 있다 \cite{guedon2024sugar}.
다만 추가 regularization/refinement 과정이 필요하고, thin open structure나 leaf-like boundary에서는 smoothing 또는 closure artifact가 생길 수 있으므로 main default보다는 alternative baseline으로 둔다.
TSDF-style extraction은 여러 view의 rendered depth를 volumetric field로 fuse한 뒤 surface를 추출하는 단순하고 renderer-agnostic한 baseline이다 \cite{curless1996volumetric}.
구현과 비교가 쉽다는 장점이 있지만, view coverage, depth quality, truncation scale에 민감하고 얇은 구조와 open boundary를 약하게 복원할 수 있으므로 fallback 또는 lower-bound baseline으로 둔다.

여기서 중요한 점은 mesh extraction 자체를 main novelty로 주장하지 않는 것이다.
이 모듈의 목적은 topology-aware wind response를 위한 internal scaffold를 만드는 것이다.
즉 최종 rendering은 여전히 3DGS로 수행하지만, simulation과 material-space correspondence를 정의하기 위해 mesh proxy를 중간 표현으로 둔다.

\subsection*{Output}

\begin{equation}
M_{\mathrm{raw}}.
\end{equation}

$M_{\mathrm{raw}}$는 아직 바로 solver에 넣기 좋은 mesh가 아니다.
Extraction 결과는 rendering surface를 따라가도록 만들어진 경우가 많아서 triangle 크기, aspect ratio, boundary loop, non-manifold artifact, thin part sampling이 simulation 관점에서는 불안정할 수 있다.
따라서 이 출력은 다음 Module 3에서 physics-quality mesh로 정리되는 중간 산출물로 본다.

\subsection*{Effect}

proxy mesh가 생기면 vanilla 3DGS에 없는 edge/face connectivity, boundary loop, local surface frame, force propagation path, artist constraint site가 생긴다. 즉, unordered Gaussian set을 connected deformable surface로 다룰 수 있게 된다.
이 단계가 없으면 wind force를 Gaussian별로 독립적으로 적용해야 하므로, 한 부분에 걸린 힘이 주변 surface로 퍼지는 cloth-like 또는 shell-like response를 만들기 어렵다.

\subsection*{Development Plan}

Mesh extraction module은 전체 pipeline을 한 번에 완성하려고 하지 않고, 독립적으로 검증 가능한 module로 개발한다.
초기 목표는 어떤 extraction method를 쓰더라도 같은 interface로 raw mesh와 품질 리포트를 내보내는 것이다.
이를 위해 먼저 공통 data contract를 정의한다.
입력은 trained 3DGS asset, camera metadata, extraction config이고, 출력은 $M_{\mathrm{raw}}$, per-vertex/per-face normal, extraction provenance, quality metrics를 포함하는 raw proxy package로 둔다.

첫 구현 단계에서는 GOF adapter를 main path로 만든다.
GOF 실행 결과를 우리 pipeline의 $M_{\mathrm{raw}}=(V_{\mathrm{raw}},E_{\mathrm{raw}},F_{\mathrm{raw}})$ 형식으로 변환하고, mesh validity, connected components, boundary loops, face aspect ratio, self-intersection 후보, Gaussian-to-surface distance histogram을 리포트한다.
이 단계의 성공 기준은 solver까지 가지 않고도 ``3DGS asset에서 추출한 mesh가 Module 3 remeshing의 입력으로 안정적으로 전달되는가''를 확인하는 것이다.

그 다음 SuGaR adapter와 TSDF adapter를 같은 interface의 comparison baseline으로 추가한다.
두 비교군은 main pipeline을 바꾸기 위한 것이 아니라, raw mesh 품질과 이후 binding/remeshing 안정성을 GOF와 비교하기 위한 ablation이다.
따라서 세 방법은 모두 동일한 downstream checker를 통과해야 하며, 최종 비교는 rendering mesh 품질만이 아니라 simulation proxy로 변환되기 쉬운지까지 함께 본다.

\section{Module 3: Simulation-Oriented Proxy Construction}

\subsection*{Input}

\begin{equation}
M_{\mathrm{raw}},\qquad \widetilde{G}^0,\qquad \mathcal{C}_{\mathrm{user}}.
\end{equation}

여기서 $M_{\mathrm{raw}}$는 extraction으로 얻은 초기 mesh이고, $\widetilde{G}^0$는 mesh가 어떤 Gaussian detail을 보존해야 하는지 알려주는 evidence source이다.
$\mathcal{C}_{\mathrm{user}}$는 pin, support, crease, rib, boundary preference처럼 사용자가 제어하고 싶은 material/constraint 정보를 담는다.
이 세 입력을 함께 쓰는 이유는 simulation mesh가 단순히 geometric surface만 따라가는 것이 아니라, 나중에 force propagation과 user control에 적합해야 하기 때문이다.

\subsection*{Operation}

raw mesh를 render-quality mesh가 아니라 physics-quality mesh로 remesh한다. 이 remeshing은 asset setup에서 한 번만 수행한다.

\begin{equation}
M_{\mathrm{raw}}\rightarrow M_{\mathrm{sim}}^0=(V^0,E,F).
\end{equation}

$M_{\mathrm{sim}}^0$는 rest state의 simulation proxy이다.
$V^0$는 canonical vertex positions이고, $E,F$는 runtime 동안 고정되는 edge/face topology이다.
중요한 점은 runtime에는 remeshing하지 않는다는 것이다.
Topology가 고정되어야 Gaussian binding, anchor index, load sample coordinate가 시간에 따라 유지된다.

simulation target edge length는 score field로 정한다.

\begin{equation}
h_{\mathrm{sim}}(p)=\clamp\left(\frac{h_0}{\sqrt{1+s_{\mathrm{sim}}(p)}},h_{\min},h_{\max}\right).
\end{equation}

$h_{\mathrm{sim}}(p)$는 위치 $p$ 주변에서 원하는 target edge length이다.
$s_{\mathrm{sim}}(p)$가 크면 보존해야 할 구조나 binding demand가 높다고 보고 edge length를 작게 만든다.
반대로 score가 낮은 완만한 영역은 더 큰 triangle을 허용하여 solver vertex 수를 줄인다.
$h_{\min}$과 $h_{\max}$는 지나치게 촘촘하거나 거친 mesh가 생기지 않도록 막는 안전한 하한/상한이다.

MVP score는 다음처럼 둘 수 있다.

\begin{equation}
s_{\mathrm{sim}}(p)=2.0B(p)+1.5K_{\mathrm{geom}}(p)+0.7R_G(p).
\end{equation}

이 score는 세 종류의 요구를 섞는다.
$B(p)$는 constraint나 boundary 보존 요구, $K_{\mathrm{geom}}(p)$는 geometry detail 보존 요구, $R_G(p)$는 Gaussian을 안정적으로 transport해야 하는 rendering demand이다.
가중치는 초기값일 뿐이며, ablation에서는 boundary-only, geometry-only, reliability-weighted demand 제거 등을 비교할 수 있다.

\subsection*{Boundary field}

\begin{equation}
B(p)=\max_{\ell\in\mathcal{L}_{\mathrm{bdry}}}w_\ell\exp\left(-\frac{d_{\mathrm{geo}}(p,\ell)^2}{2r_B^2}\right).
\end{equation}

이 항은 flag pole edge, curtain rail, leaf stem, pinned boundary 같은 constraint-critical region을 보존한다.
$\mathcal{L}_{\mathrm{bdry}}$는 사용자가 지정하거나 자동 검출한 boundary/support curve들의 집합이다.
$d_{\mathrm{geo}}(p,\ell)$는 surface 위에서 point $p$와 boundary curve $\ell$ 사이의 geodesic distance이고, $r_B$는 boundary 영향이 퍼지는 폭이다.
Euclidean distance가 아니라 geodesic distance를 쓰는 이유는 얇은 접힘이나 가까운 반대편 surface가 잘못 같은 boundary 영향권으로 묶이는 것을 줄이기 위해서이다.

\subsection*{Geometry field}

face normal variation 기반으로 계산할 수 있다.

\begin{equation}
K_{\mathrm{geom}}(f)=\frac{1}{|\mathcal{N}(f)|}\sum_{f'\in\mathcal{N}(f)}\frac{1-\langle n_f,n_{f'}\rangle}{\|c_f-c_{f'}\|+\eps}.
\end{equation}

$K_{\mathrm{geom}}(f)$는 face $f$ 주변에서 normal이 얼마나 빠르게 변하는지 보는 curvature-like score이다.
$\mathcal{N}(f)$는 neighboring faces, $n_f$와 $c_f$는 각각 face normal과 centroid이다.
두 face normal이 비슷하면 $1-\langle n_f,n_{f'}\rangle$가 작고, fold나 sharp edge처럼 normal이 급격히 바뀌면 값이 커진다.
거리로 나누는 이유는 같은 normal 변화라도 짧은 거리에서 발생하면 더 날카로운 구조로 보기 위해서이다.

robust normalization은 다음처럼 둔다.

\begin{equation}
\widehat{K}_{\mathrm{geom}}(p)=\clamp\left(\frac{K_{\mathrm{geom}}(p)-P_5}{P_{95}-P_5+\eps},0,1\right).
\end{equation}

이 항은 leaf tip, cloth fold, paper edge, narrow strip 같은 thin or high-curvature structure를 보존한다.
$P_5$와 $P_{95}$를 쓰는 이유는 소수의 outlier curvature가 전체 normalization을 망치는 것을 막기 위해서이다.
실제 remeshing score에는 normalized $\widehat{K}_{\mathrm{geom}}$를 넣는 것이 안전하며, 본 문서에서 $K_{\mathrm{geom}}$라고 쓸 때는 필요에 따라 normalized field를 의미하도록 약속할 수 있다.

\subsection*{Reliability-weighted Gaussian transport demand}

raw Gaussian density 대신, surface-consistent하고 transport할 가치가 있는 Gaussian의 demand를 사용한다.

\begin{equation}
R_G(p)=1-\exp\left(-\frac{\sum_j q_j\eta_jK_j(p)}{\tau_R}\right).
\end{equation}

\begin{equation}
K_j(p)=\exp\left(-\frac{\|p-p_j\|^2}{2r_j^2}\right),
\end{equation}

여기서 $p_j$는 Gaussian center $\mu_j^0$를 proxy surface에 projection한 point이고, $r_j$는 surface footprint radius이다. $\eta_j$는 rendering contribution이며, MVP에서는 $\eta_j=1$로 둘 수 있다. $\tau_R$는 saturation parameter로 $P_{90}$ 같은 percentile을 사용할 수 있다.
이 항의 의도는 단순히 Gaussian이 많은 곳을 중요하다고 보는 것이 아니다.
Gaussian 수는 reconstruction 품질, optimizer behavior, noisy floater에 영향을 받을 수 있으므로 raw density만 쓰면 잘못된 곳을 과하게 보존할 수 있다.
따라서 $q_j$로 reliability를 반영하고, $\eta_j$로 실제 rendering contribution을 반영한 뒤, $K_j(p)$로 surface 위치 $p$ 주변에 영향을 퍼뜨린다.
$1-\exp(\cdot)$ 형태는 demand가 일정 수준 이상 커지면 saturation되도록 하여 한 영역이 과도하게 remesh되는 것을 막는다.

\subsection*{Gated floater likelihood}

isolation은 직접 reliability factor로 쓰지 않는다. 대신 low opacity, far from surface, isolated, and not protected인 경우에만 floater likelihood로 작동한다.

\begin{align}
q_j &= q_j^{\mathrm{base}}\exp(-\lambda_{\mathrm{float}}p_j^{\mathrm{float}}),\\
p_j^{\mathrm{float}} &= I_jL_jD_j(1-S_j),\\
I_j &= \exp\left(-\frac{N_j^{\mathrm{nbr}}}{N_{\mathrm{ref}}}\right),\\
L_j &= 1-q_j^\alpha,\\
D_j &= 1-q_j^{\mathrm{surf}},\\
q_j^{\mathrm{surf}} &= \exp\left(-\frac{d_j^2}{2\sigma_d^2}\right).
\end{align}

$q_j^{\mathrm{base}}$는 opacity, view contribution, surface consistency 등에서 얻은 기본 reliability이다.
$p_j^{\mathrm{float}}$는 floater일 가능성을 나타내는 soft penalty이고, 네 조건이 동시에 맞을 때만 커진다.
$I_j$는 주변 Gaussian이 적을수록 커지는 isolation score, $L_j$는 low-opacity일수록 커지는 score, $D_j$는 surface에서 멀수록 커지는 score이다.
마지막 $(1-S_j)$는 보호되어야 하는 구조에서는 penalty를 끄기 위한 gate이다.
이 곱 구조 덕분에 ``isolated but valid'' thin detail을 단순히 제거하지 않는다.

protection score는 다음처럼 정의한다.

\begin{equation}
S_j=\max_{s\in\mathcal{S}_j}s,
\end{equation}

\begin{equation}
\begin{aligned}
\mathcal{S}_j={}&
\{B(p_j),K_{\mathrm{geom}}(p_j),C_{\mathrm{constraint}}(p_j)\}\\
&\cup\{T_{\mathrm{thin}}(p_j)\;\text{if implemented}\}\\
&\cup\{q_j^{\mathrm{view}}\;\text{if measured or approximated}\}.
\end{aligned}
\end{equation}

주의할 점은 $q_j^{\mathrm{view}}$가 없을 때 base reliability에서는 neutral factor로 1을 둘 수 있지만, $S_j$의 max에는 넣지 않는다는 것이다. 그렇지 않으면 $S_j=1$이 되어 floater penalty가 항상 꺼진다.
$S_j$는 boundary, high-curvature geometry, user constraint, thin-structure detector, measured view contribution 중 하나라도 강하면 Gaussian을 보호한다.
따라서 floater filtering은 removal rule이 아니라 risk weighting rule이다.
이 관점은 valid sparse detail과 noise를 구분해야 하는 3DGS preprocessing에서 특히 중요하다.

\subsection*{Output}

\begin{equation}
M_{\mathrm{sim}}^0=(V^0,E,F).
\end{equation}

$M_{\mathrm{sim}}^0$는 이후 모든 runtime module이 공유하는 fixed-topology proxy이다.
이 mesh의 vertex/face index는 simulation anchors, load samples, Gaussian bindings의 reference index가 되므로 asset setup 이후에는 바뀌지 않는다.

\subsection*{Effect}

이 모듈은 runtime remeshing 없이 사용할 compact simulation mesh를 만든다. 결과적으로 boundary, thin structure, Gaussian binding quality, simulation stability를 동시에 확보한다.
렌더링 품질만 보면 원본 Gaussian이 담당하고, dynamics 품질은 $M_{\mathrm{sim}}^0$가 담당한다.
따라서 이 모듈의 성공 기준은 mesh가 시각적으로 가장 예쁜지가 아니라, stable solver, stable binding, controllable constraint를 동시에 만족하는지이다.

\section{Module 4: Simulation Anchors and Load Samples}

\subsection*{Input}

\begin{equation}
M_{\mathrm{sim}}^0=(V^0,E,F).
\end{equation}

이 모듈은 이미 topology가 정리된 simulation proxy를 입력으로 받는다.
목표는 같은 mesh 위에서 세 종류의 고정 material-space reference를 정의하는 것이다: solver가 움직일 simulation anchors, 사용자가 고정하거나 제어할 constraint anchors, wind load를 샘플링할 load samples이다.

\subsection*{Operation: simulation anchors}

simulation anchor는 mesh vertex이다.

\begin{equation}
A_{\mathrm{sim}}=V.
\end{equation}

$A_{\mathrm{sim}}$는 solver가 실제로 상태를 적분하는 자유도 집합이다.
초기 MVP에서는 proxy mesh vertex 전체를 simulation anchor로 쓰는 것이 가장 단순하고 안정적이다.
나중에 더 coarse한 anchor graph를 쓰고 싶다면 $A_{\mathrm{sim}}\subset V$로 줄일 수 있지만, 그러면 interpolation과 force projection이 추가로 필요하다.

각 anchor state는 다음과 같다.

\begin{equation}
a_i(t)=(x_i(t),v_i(t),m_i,\theta_i).
\end{equation}

$x_i(t)$와 $v_i(t)$는 runtime position/velocity이고, $m_i$는 vertex에 할당된 mass이다.
$\theta_i$는 stiffness, damping, drag response, density 같은 material parameter를 담는 자리이다.
이 state는 Gaussian state가 아니라 proxy simulation state이다.
Gaussian은 직접 물리 solver의 자유도가 되지 않고, 뒤의 transport module에서 이 proxy state를 따라간다.

area와 mass는 incident faces로부터 계산한다.

\begin{equation}
A_i=\sum_{T\ni i}\frac{\operatorname{area}(T)}{3},\qquad m_i=\rho A_i.
\end{equation}

$A_i$는 vertex $i$가 대표하는 surface area이다.
각 triangle area를 세 vertex에 균등 분배하면 간단한 lumped mass를 만들 수 있다.
$\rho$는 areal density이며, cloth, paper, rubber-like shell 등 profile preset에 따라 바꿀 수 있다.

material parameter는 초기에는 per-vertex field보다 material-class scalar로 시작하는 편이 안전하다.

\begin{equation}
\theta_i\approx\theta_{\mathrm{mat}}=\{\alpha,d,\beta,\rho\}.
\end{equation}

여기서 $\alpha$는 compliance 또는 inverse stiffness, $d$는 damping, $\beta$는 bending-related parameter, $\rho$는 density로 둘 수 있다.
초기에는 모든 vertex에 같은 material-class scalar를 부여하고, 이후 controllable material field가 안정되면 per-vertex 또는 painted field로 확장한다.

\subsection*{Operation: constraint anchors}

사용자 pin 또는 support는 subset이다.

\begin{equation}
A_{\mathrm{pin}}\subset A_{\mathrm{sim}}.
\end{equation}

$A_{\mathrm{pin}}$은 solver 자유도 중 움직임을 제한하거나 target에 붙일 anchor subset이다.
예를 들어 연의 뼈대, 종이비행기의 접힌선, 깃발의 pole edge, leaf stem 같은 부분은 사용자 constraint나 profile preset으로 지정할 수 있다.
이 constraint anchor는 앞서 말한 simulation anchor와 목적이 다르다.
Simulation anchor는 solver의 계산 자유도이고, constraint anchor는 그 자유도 중 제어 조건이 걸리는 subset이다.

hard pin은:

\begin{equation}
x_i(t)=x_i^0,\qquad i\in A_{\mathrm{pin}}.
\end{equation}

Hard pin은 해당 vertex가 rest position에서 움직이지 않는다는 의미이다.
완전히 고정된 support나 boundary를 표현할 때 좋지만, 너무 많이 쓰면 deformation이 인위적으로 보일 수 있다.

soft pin은:

\begin{equation}
C_i^{\mathrm{pin}}(x)=x_i-x_i^{\mathrm{target}}.
\end{equation}

Soft pin은 target과의 차이를 constraint로 두고 compliance를 통해 어느 정도 따라가게 하는 방식이다.
사용자가 handle을 움직이거나, 뼈대가 완전 고정이 아니라 약간 휘는 경우에는 hard pin보다 soft pin이 더 자연스럽다.

\subsection*{Operation: fixed material-space load samples}

wind force는 반드시 simulation vertex에서만 계산하지 않는다. face 위에 fixed material-space load sample을 둔다.

\begin{equation}
q_l=(T_l,b_l).
\end{equation}

$q_l$는 $l$번째 load sample의 material-space coordinate이다.
$T_l$은 sample이 놓인 rest-state triangle index이고, $b_l$은 그 triangle 내부의 barycentric coordinate이다.
따라서 sample은 월드 좌표에 고정된 점이 아니라, mesh material에 붙어 있는 점이다.
바람에 의해 mesh가 움직여도 같은 삼각형 위의 같은 재료 위치를 계속 따라간다.

runtime 위치는:

\begin{equation}
x_l(t)=\sum_{i\in T_l}b_{li}x_i(t).
\end{equation}

이 식은 barycentric interpolation으로 load sample의 현재 위치를 계산한다.
중요한 점은 exposure나 wind가 frame마다 바뀌어도 $q_l=(T_l,b_l)$ 자체는 바뀌지 않는다는 것이다.
즉 우리는 runtime adaptive resampling이 아니라 fixed material-space sampling을 사용한다.
이 선택 덕분에 force history와 temporal correspondence가 안정적으로 유지된다.

\subsection*{Output}

\begin{equation}
A_{\mathrm{sim}},\qquad A_{\mathrm{pin}},\qquad Q_{\mathrm{load}}=\{q_l\}.
\end{equation}

$A_{\mathrm{sim}}$는 solver 자유도, $A_{\mathrm{pin}}$은 constraint가 걸리는 자유도 subset, $Q_{\mathrm{load}}$는 wind force를 평가할 material-space samples이다.
세 집합은 모두 rest proxy mesh 위에서 한 번 정의되고 runtime 동안 유지된다.

\subsection*{Effect}

이 모듈은 simulation resolution과 force sampling resolution을 분리한다. 따라서 wind load가 복잡한 곳에는 load sample만 촘촘히 둘 수 있고, XPBD solver vertex 수는 작게 유지할 수 있다.
이 구분이 본래 anchor 설계 의도와 연결된다.
모든 지점에서 비싼 force evaluation이나 dense simulation을 하지 않고, 중요한 material-space sample과 compact solver anchors 사이를 projection으로 연결한다.

\section{Module 5: Gaussian-to-Proxy Binding}

\subsection*{Input}

\begin{equation}
G^0,\qquad M_{\mathrm{sim}}^0.
\end{equation}

입력 $G^0$는 rendering primitive인 Gaussian set이고, $M_{\mathrm{sim}}^0$는 dynamics를 담당하는 proxy mesh이다.
이 모듈은 두 representation 사이의 persistent correspondence를 만든다.
한 번 binding이 만들어지면 runtime에는 Gaussian이 매 frame 새 triangle을 찾지 않고, 저장된 material coordinate를 따라간다.

\subsection*{Operation}

각 Gaussian을 가장 안정적인 mesh triangle에 binding한다.

\begin{equation}
B(j)=\argmin_{T\in F}E_{jT}.
\end{equation}

$B(j)$는 Gaussian $j$가 어떤 triangle $T$에 붙을지를 결정하는 mapping이다.
가장 가까운 triangle만 고르면 surface 반대편이나 noisy mesh에 잘못 붙을 수 있으므로, distance뿐 아니라 normal consistency와 triangle quality를 함께 본다.

binding energy 예시는 다음과 같다.

\begin{equation}
E_{jT}=\frac{d(\mu_j^0,T)^2}{\sigma_j^2}+\lambda_n(1-|\langle n_j^G,n_T\rangle|)+\lambda_qQ_T.
\end{equation}

$d(\mu_j^0,T)$는 Gaussian center와 triangle 사이의 거리이고, $\sigma_j$는 Gaussian scale 기반 normalization이다.
두 번째 항은 Gaussian normal candidate $n_j^G$와 triangle normal $n_T$의 정렬도를 본다.
절댓값을 쓰는 이유는 covariance에서 얻은 eigenvector normal의 부호가 모호하기 때문이다.
$Q_T$는 skinny triangle, boundary artifact, low-quality face에 대한 penalty로 둘 수 있다.

각 Gaussian에 대해 저장한다.

\begin{equation}
\mathcal{B}_j=(T_j,b_j,R_{T_j}^0,\delta_j).
\end{equation}

여기서 $T_j=(a,b,c)$, $b_j=(b_{ja},b_{jb},b_{jc})$는 barycentric coordinate, $R_{T_j}^0$는 canonical triangle local frame, $\delta_j$는 optional local offset이다.
$b_j$는 Gaussian center가 triangle 위의 어느 material position에 붙는지를 나타내고, $R_{T_j}^0$는 triangle의 rest local frame이다.
$\delta_j$를 두면 Gaussian이 surface 위에 정확히 놓이지 않고 약간 떠 있거나 두께 방향 offset을 갖는 경우도 표현할 수 있다.
이 정보가 있어야 뒤에서 deformed triangle frame을 이용해 mean과 covariance를 함께 transport할 수 있다.

\subsection*{Output}

\begin{equation}
\{\mathcal{B}_j\}_{j=1}^{N_G}.
\end{equation}

출력은 Gaussian마다 하나씩 저장된 binding table이다.
이 table은 runtime 동안 고정되며, Gaussian index $j$와 proxy triangle/material coordinate 사이의 correspondence를 보장한다.

\subsection*{Effect}

persistent binding은 mesh deformation을 Gaussian rendering representation으로 되돌리기 위한 핵심이다. 이 단계가 있어야 per-frame closest triangle search나 rebinding 없이 temporal consistency를 유지할 수 있다. 또한 Gaussian mean, anisotropic covariance, local frame, SH feature를 mesh deformation에 따라 일관되게 transport할 수 있다.
논문 관점에서는 이 부분이 ``3DGS와 simulation proxy 사이의 bridge''를 이루는 핵심 구성요소이다.
Binding이 불안정하면 solver가 아무리 좋아도 rendering 단계에서 Gaussian이 튀거나 surface에서 미끄러지는 artifact가 생긴다.

\section{Module 6: Blender Dense Supervision Generation}

\subsection*{Input}

Blender 또는 physics engine의 known mesh, material, wind condition, camera views.

\begin{equation}
M_{\mathrm{hi}}^0,\qquad W(x,t),\qquad \Pi_t.
\end{equation}

$M_{\mathrm{hi}}^0$는 high-resolution reference mesh 또는 simulator-native asset이다.
이 모듈의 입력은 실제 target 3DGS asset이라기보다, 학습과 평가를 위해 ground-truth force/deformation/rendering을 얻을 수 있는 synthetic setup이다.
$W(x,t)$는 simulator에 넣는 wind condition이고, $\Pi_t$는 rendered reference를 만들 camera sequence이다.

\subsection*{Operation}

각 asset에 대해 여러 wind condition을 simulation하고, dense state와 rendered reference를 저장한다.

\begin{equation}
\mathcal{D}_{\mathrm{Blender}}=\{x_r^{\mathrm{dense}}(t),f_r^{\mathrm{dense}}(t),W(t),I_t^{\mathrm{gt}}\}.
\end{equation}

$x_r^{\mathrm{dense}}(t)$는 dense mesh vertex 또는 sampled point의 trajectory이고, $f_r^{\mathrm{dense}}(t)$는 같은 위치에서의 reference force이다.
$W(t)$는 해당 sequence의 wind parameter 또는 field sample이고, $I_t^{\mathrm{gt}}$는 camera $\Pi_t$에서 렌더링한 reference image이다.
이 데이터는 runtime에 직접 쓰기 위한 것이 아니라, force predictor와 projection module을 검증하거나 학습하기 위한 supervision이다.

\subsection*{Output}

\begin{equation}
\mathcal{D}_{\mathrm{Blender}}.
\end{equation}

출력 dataset은 dense physical signal과 image-space signal을 함께 갖는다.
Dense force는 load prediction/projection의 target으로 쓰고, rendered image는 최종 visual response가 reference와 얼마나 맞는지 평가하는 데 쓴다.

\subsection*{Effect}

Blender supervision은 runtime engine이 아니라 학습 데이터와 evaluation ground truth를 제공한다. 이 데이터는 force predictor, dense-to-proxy projection, SH residual, ablation evaluation에 사용된다.
즉 이 모듈은 논문 시스템의 runtime dependency가 아니다.
Asset이 실제 captured 3DGS뿐인 경우에는 이 모듈을 건너뛰고 analytic force baseline부터 사용할 수 있으며, synthetic benchmark를 만들 때만 dense supervision을 생성한다.

\section{Module 7: Conservation-Preserving Dense-to-Proxy Force Projection}

\subsection*{Input}

Blender dense force와 proxy/load sample layout.

\begin{equation}
\{f_r^w\}_{r\in\Omega},\qquad \{x_i\}_{i\in\Omega}.
\end{equation}

$\Omega$는 dense reference force를 하나의 proxy patch 또는 load-sampling region으로 모은 영역이다.
$f_r^w$는 wind condition $w$에서 dense sample $r$에 작용한 reference force이고, $x_i$는 그 patch를 대표하는 proxy/load sample positions이다.
이 모듈의 목적은 dense supervision을 sparse proxy representation으로 줄이되, 전체 force와 torque 정보를 잃지 않는 것이다.

\subsection*{Operation}

patch $\Omega$에서 dense net force와 wrench를 계산한다.

\begin{align}
F_{\Omega,w}^{\mathrm{dense}} &= \sum_{r\in\Omega}f_r^w,\\
T_{\Omega,w}^{\mathrm{dense}} &= \sum_{r\in\Omega}(x_r-c_\Omega)\times f_r^w.
\end{align}

$F_{\Omega,w}^{\mathrm{dense}}$는 patch 전체에 작용하는 net force이고, $T_{\Omega,w}^{\mathrm{dense}}$는 patch center $c_\Omega$ 기준의 torque 또는 wrench moment이다.
단순히 force magnitude만 맞추면 patch가 회전하려는 경향을 잃을 수 있으므로, torque 보존도 함께 둔다.

raw proxy force $F_{i,w}^{\mathrm{raw}}$를 얻은 뒤, conservation-preserving correction을 푼다.

\begin{equation}
\min_{\{F_{i,w}^{*},\tau_{i,w}^{*}\}}\sum_i\|F_{i,w}^{*}-F_{i,w}^{\mathrm{raw}}\|^2+\gamma\|\tau_{i,w}^{*}-\tau_{i,w}^{\mathrm{raw}}\|^2,
\end{equation}

subject to:

\begin{align}
\sum_iF_{i,w}^{*} &= F_{\Omega,w}^{\mathrm{dense}},\\
\sum_i(x_i-c_\Omega)\times F_{i,w}^{*}+\tau_{i,w}^{*} &= T_{\Omega,w}^{\mathrm{dense}}.
\end{align}

$F_{i,w}^{\mathrm{raw}}$는 nearest assignment, kernel interpolation, or coarse predictor로 얻은 초기 sparse force이다.
최적화는 이 raw force를 가능한 적게 바꾸면서, patch 전체의 dense net force와 torque가 정확히 맞도록 보정한다.
$\tau_{i,w}^{*}$는 force만으로 torque를 완전히 맞추기 어려울 때 남는 local torque term으로 볼 수 있다.
이 step은 sparse representation이 physical aggregate를 보존한다는 근거가 된다.

projection risk는 initial coarse load-sampling layout 또는 patch partition에서 계산한 correction ratio로 정의한다.

\begin{equation}
\epsilon_{\Omega,w}^{\mathrm{proj}}=\frac{\sum_i\|F_{i,w}^{*}-F_{i,w}^{\mathrm{raw}}\|^2+\gamma\sum_i\|\tau_{i,w}^{*}-\tau_{i,w}^{\mathrm{raw}}\|^2}{\sum_i\|F_{i,w}^{\mathrm{raw}}\|^2+\gamma\sum_i\|\tau_{i,w}^{\mathrm{raw}}\|^2+\eps}.
\end{equation}

$\epsilon_{\Omega,w}^{\mathrm{proj}}$가 크다는 것은 coarse layout이 dense force를 잘 표현하지 못해 correction이 많이 필요했다는 뜻이다.
따라서 이 값은 어디에 load sample을 더 두어야 하는지 알려주는 diagnostic signal이다.
분모의 $\eps$는 raw force가 거의 0인 영역에서 ratio가 폭주하는 것을 막는다.

\begin{equation}
F_{\mathrm{projection\_risk}}(p)=\normalize\left(\max_{w\in\mathcal{W}}\epsilon_{\Omega(p),w}^{\mathrm{proj}}\right).
\end{equation}

이 값은 runtime adaptive sampling이 아니다. Preprocessing에서 coarse layout 위에 한 번 계산하고, fixed load sample placement를 한 번 refine하는 데 사용한다.
$\max_{w\in\mathcal{W}}$를 쓰는 이유는 여러 wind condition 중 하나라도 표현하기 어려운 영역이면 load sample refinement 후보로 잡기 위해서이다.
즉 projection risk는 time-varying anchor가 아니라 fixed layout을 더 잘 설계하기 위한 one-time preprocessing score이다.

\begin{equation}
s_{\mathrm{load}}(p)=w_EE_{\mathrm{prior}}(p)+w_F\,F_{\mathrm{projection\_risk}}(p)+w_RR_G(p)+w_KK_{\mathrm{geom}}(p).
\end{equation}

$s_{\mathrm{load}}(p)$는 load sample density를 정하는 score이다.
$E_{\mathrm{prior}}$는 wind exposure나 silhouette-like prior, $F_{\mathrm{projection\_risk}}$는 dense force 보존 난이도, $R_G$는 Gaussian transport/rendering demand, $K_{\mathrm{geom}}$는 geometry detail demand를 나타낸다.
이 score는 solver mesh resolution과 별개로 load sample placement를 조절한다.

Blender-supervised MVP에서는 다음처럼 시작할 수 있다.

\begin{equation}
s_{\mathrm{load}}^{\mathrm{Blender\text{-}MVP}}(p)=1.2F_{\mathrm{projection\_risk}}(p)+0.5K_{\mathrm{geom}}(p).
\end{equation}

초기에는 Blender에서 얻은 projection risk와 geometry curvature만으로도 어느 영역의 load sampling이 부족한지 확인할 수 있다.
Blender supervision이 없는 실제 captured asset에서는 analytic exposure prior와 geometry/Gaussian demand로 대체한다.

\subsection*{Output}

\begin{equation}
F_i^*,\quad \tau_i^*,\quad F_{\mathrm{projection\_risk}}(p),\quad Q_{\mathrm{load}}^{\mathrm{refined}}.
\end{equation}

출력은 보정된 sparse force/torque와, 그 보정량으로부터 얻은 projection risk field, 그리고 한 번 refine된 fixed load sample set이다.
$Q_{\mathrm{load}}^{\mathrm{refined}}$도 runtime에는 고정된다.

\subsection*{Effect}

이 모듈은 dense Blender force를 coarse proxy/load samples로 압축하면서 net force와 wrench를 보존한다. 따라서 sparse proxy simulation이 단순 계산 shortcut이 아니라, dense force supervision을 보존하는 representation으로 정당화된다.
논문 실험에서는 이 모듈을 통해 random/coarse sampling 대비 projection-risk-aware load sampling이 force preservation과 visual response를 개선하는지 ablation할 수 있다.

\section{Module 8: Runtime Wind Field Evaluation}

\subsection*{Input}

\begin{equation}
M(t),\qquad W(x,t),\qquad Q_{\mathrm{load}}.
\end{equation}

$M(t)$는 현재 frame의 deformed proxy mesh이고, $W(x,t)$는 사용자나 procedural controller가 정의한 wind velocity field이다.
$Q_{\mathrm{load}}$는 Module 4 또는 7에서 만든 fixed material-space load samples이다.
이 모듈은 mesh topology나 load sample placement를 바꾸지 않고, 현재 위치에서 wind를 평가하는 runtime query 단계이다.

\subsection*{Operation}

wind field는 lightweight analytic field로 둔다.

\begin{equation}
\begin{aligned}
W(x,t)=&\,W_{\mathrm{global}}(t)
+\sum_kW_{\mathrm{source},k}(x,t)\\
&+W_{\mathrm{gust}}(x,t).
\end{aligned}
\end{equation}

$W_{\mathrm{global}}(t)$는 전체 장면에 적용되는 direction/strength이며, fan이나 directional wind control로 볼 수 있다.
$W_{\mathrm{source},k}(x,t)$는 위치 의존적인 local source, 예를 들어 fan, vortex, brush stroke wind를 표현한다.
$W_{\mathrm{gust}}(x,t)$는 sinusoidal gust나 turbulence noise처럼 시간적으로 흔들리는 성분이다.
이 설계는 full CFD를 복원하려는 것이 아니라, controllable one-way wind input을 제공하는 데 목적이 있다.

각 load sample 위치와 velocity를 계산한다.

\begin{align}
x_l(t)&=\sum_{i\in T_l}b_{li}x_i(t),\\
v_l(t)&=\sum_{i\in T_l}b_{li}v_i(t),\\
W_l(t)&=W(x_l(t),t).
\end{align}

$x_l(t)$와 $v_l(t)$는 load sample $q_l=(T_l,b_l)$이 현재 deformed mesh에서 어디에 있고 얼마나 움직이는지 나타낸다.
둘 다 barycentric interpolation으로 계산하므로, sample은 같은 material point를 계속 따라간다.
$W_l(t)$는 그 현재 위치에서 wind field를 query한 값이다.
이 단계까지는 wind와 object motion을 별도로 계산하며, 실제 상대 바람은 다음 exposure module에서 만든다.

\subsection*{Output}

\begin{equation}
W_l(t),\qquad x_l(t),\qquad v_l(t).
\end{equation}

출력은 load sample별 current position, current material velocity, local wind velocity이다.
이 세 값은 aerodynamic force를 계산하기 위한 최소 runtime feature이다.

\subsection*{Effect}

user wind input을 proxy surface의 local wind로 변환한다. Full CFD를 주장하지 않으면서 controllable wind editing을 가능하게 한다.
사용자가 wind direction, gust frequency, fan source를 조절하면 이 모듈의 $W_l(t)$가 바뀌고, 이후 force와 deformation response가 달라진다.

\section{Module 9: Wind Exposure Estimation}

\subsection*{Input}

\begin{equation}
G(t)\;\text{or}\;M(t),\qquad \hat{w}(t),\qquad Q_{\mathrm{load}}.
\end{equation}

입력은 현재 deformed Gaussian state 또는 proxy mesh, wind direction $\hat{w}(t)$, 그리고 fixed load samples이다.
Exposure estimation은 단순히 wind magnitude를 쓰는 것이 아니라, 바람 방향에서 보았을 때 각 load sample이 얼마나 직접적으로 바람을 맞는지 추정한다.
front surface, occluded surface, backside surface가 같은 force를 받지 않도록 하기 위한 단계이다.

\subsection*{Operation}

current Gaussian or proxy state를 wind-axis coordinate에서 rasterize하여 transmittance map을 만든다.

\begin{equation}
S_t(u)=\exp\left(-\int\alpha_t(s)ds\right).
\end{equation}

$u$는 wind direction에 수직인 2D wind-axis image coordinate이고, $s$는 wind ray를 따라가는 depth coordinate이다.
$\alpha_t(s)$는 그 ray 위에서 만나는 opacity density이다.
$S_t(u)$는 바람이 해당 ray를 따라 들어올 때 얼마나 통과하는지를 나타내는 transmittance이다.
값이 1에 가까우면 wind-facing exposure가 크고, 0에 가까우면 앞쪽 surface에 의해 막혀 있다고 본다.

Gaussian splatting approximation은 다음과 같다.

\begin{equation}
S_t(u)\approx\prod_j(1-\alpha_jK_j^w(u)).
\end{equation}

$K_j^w(u)$는 Gaussian $j$를 wind-axis plane에 projection했을 때의 footprint kernel이다.
이 approximation은 3DGS의 alpha compositing 직관을 wind direction으로 돌려 쓴 것이다.
정확한 fluid occlusion은 아니지만, self-shadowing과 wind blocking을 가볍게 feature화할 수 있다.

각 fixed load sample에서 exposure를 sampling한다.

\begin{equation}
e_l(t)=\operatorname{sample}(S_t,\pi_w(x_l(t))).
\end{equation}

$\pi_w$는 3D point를 wind-axis plane으로 projection하는 함수이다.
따라서 $e_l(t)$는 load sample $l$이 현재 frame에서 얼마나 wind-exposed인지 나타내는 scalar이다.
중요하게도 sample 위치 $x_l(t)$는 움직이지만 sample identity $q_l$은 고정되어 있다.

relative effective wind는:

\begin{equation}
u_l^{\mathrm{eff}}(t)=e_l(t)W_l(t)-v_l(t).
\end{equation}

$u_l^{\mathrm{eff}}(t)$는 object material point가 실제로 느끼는 상대 바람이다.
$e_l(t)W_l(t)$는 exposure로 감쇠된 wind velocity이고, $v_l(t)$는 material point 자체의 velocity이다.
물체가 wind 방향으로 같이 움직이면 상대 wind가 줄고, 반대로 움직이면 더 큰 force가 생긴다.

\subsection*{Output}

\begin{equation}
e_l(t),\qquad u_l^{\mathrm{eff}}(t).
\end{equation}

출력은 force predictor가 바로 사용할 exposure scalar와 relative wind vector이다.
$e_l(t)$는 load magnitude뿐 아니라 SH residual이나 appearance correction의 feature로도 사용할 수 있다.

\subsection*{Effect}

exposure는 self-shadowing, wind blocking, front/back visibility를 force feature로 반영한다. 중요한 점은 exposure가 frame마다 변해도 anchor나 load sample을 재배치하지 않는다는 것이다. Fixed material-space sample에서 값만 갱신한다.
이 설계는 사용자가 제기했던 ``wind field가 frame마다 바뀌면 anchor도 바뀌어야 하는가''라는 문제를 피한다.
Anchor와 load sample은 고정하고, 그 위에서 exposure value만 time-varying field로 평가한다.

\section{Module 10: Force-Space Load Prediction}

\subsection*{Input}

\begin{equation}
z_l(t),\qquad u_l^{\mathrm{eff}}(t),\qquad e_l(t),\qquad \theta_l.
\end{equation}

$z_l(t)$는 load sample 주변의 local geometry/state feature이다.
예를 들어 normal, area, curvature, strain, previous velocity, local frame 등을 포함할 수 있다.
$u_l^{\mathrm{eff}}(t)$와 $e_l(t)$는 wind-related feature이고, $\theta_l$은 material parameter 또는 profile preset에서 온 local material feature이다.

\subsection*{Operation}

deformation state를 직접 예측하지 않고 aerodynamic load를 예측한다.

\begin{equation}
(f_l(t),\tau_l(t))=P_\theta(z_l(t),u_l^{\mathrm{eff}}(t),e_l(t),\theta_l).
\end{equation}

$P_\theta$는 analytic formula, learned network, or analytic-plus-residual model일 수 있다.
출력 $f_l(t)$는 load sample에 작용하는 force이고, $\tau_l(t)$는 필요할 경우 local torque 또는 bending moment이다.
우리가 deformation을 직접 regression하지 않는 이유는 solver가 topology와 constraints를 처리하도록 남겨두기 위해서이다.

analytic drag baseline은 다음처럼 둘 수 있다.

\begin{equation}
\begin{aligned}
f_l^{\mathrm{drag}}(t)=&
\frac{1}{2}\rho_{\mathrm{air}}C_DA_l e_l(t)
\|u_l^{\mathrm{eff}}(t)\|^2\\
&\cdot \max(0,\langle n_l(t),\hat{u}_l(t)\rangle)n_l(t).
\end{aligned}
\end{equation}

$\rho_{\mathrm{air}}$는 air density, $C_D$는 drag coefficient, $A_l$은 load sample이 대표하는 area이다.
$\max(0,\langle n_l,\hat{u}_l\rangle)$는 wind-facing surface만 force를 받도록 하는 orientation gate이다.
이 baseline은 물리적으로 완전한 aerodynamics는 아니지만, wind strength가 제곱에 비례하고 surface orientation에 의존한다는 기본 경향을 제공한다.

학습형 force predictor는 residual 형태가 안전하다.

\begin{equation}
f_l(t)=f_l^{\mathrm{analytic}}(t)+\Delta f_l^\theta(t).
\end{equation}

Residual 형태는 analytic baseline이 설명하지 못하는 flutter, vortex-like response, material-specific correction만 학습하게 한다.
이렇게 하면 network가 전체 dynamics를 외우는 것보다 wind control과 out-of-distribution amplitude에 더 안정적일 가능성이 높다.

\subsection*{Output}

\begin{equation}
f_l(t),\qquad \tau_l(t).
\end{equation}

출력은 아직 solver vertex force가 아니라 load sample force이다.
다음 Module 11에서 barycentric projection을 통해 simulation anchors로 누적된다.

\subsection*{Effect}

이 모듈은 direct deformation regression을 피하고, wind amplitude/direction/frequency에 대해 더 제어 가능한 response를 만든다. 또한 Blender dense force supervision과 자연스럽게 연결된다.
논문 framing에서는 ``learned deformation''보다 ``force-space wind response prior''에 가깝게 설명하는 편이 안전하다.

\section{Module 11: Load-to-Anchor Force Projection}

\subsection*{Input}

\begin{equation}
f_l(t),\quad \tau_l(t),\quad q_l=(T_l,b_l).
\end{equation}

입력 force는 load sample 위치에서 계산된 값이고, XPBD solver는 vertex anchor force를 필요로 한다.
따라서 이 모듈은 material-space load sample과 simulation anchor 사이의 force transfer를 담당한다.
$q_l=(T_l,b_l)$가 있기 때문에 sample force를 어떤 triangle의 어떤 vertex들로 보낼지 알 수 있다.

\subsection*{Operation}

기본적으로 barycentric distribution을 사용한다.

\begin{equation}
f_i(t)\mathrel{+}=b_{li}f_l(t),\qquad i\in T_l.
\end{equation}

이 식은 load sample force를 triangle vertices에 barycentric weight로 분배한다.
Sample이 vertex에 가까우면 해당 vertex가 더 큰 force를 받고, triangle 중앙이면 세 vertex가 비슷하게 나눠 받는다.
이 방식은 간단하지만 material coordinate와 잘 맞고, total force 보존도 직관적이다.

external force는 다음과 같다.

\begin{equation}
f_i^{\mathrm{ext}}(t)=\sum_{l:i\in T_l}b_{li}f_l(t)+f_i^{\mathrm{gravity}}+f_i^{\mathrm{user}}.
\end{equation}

$f_i^{\mathrm{ext}}(t)$는 solver가 다음 step을 예측할 때 사용하는 외력이다.
wind load 외에도 gravity, user handle, scripted interaction force를 같은 형태로 합산할 수 있다.
Torque $\tau_l(t)$를 사용하는 경우에는 equivalent force couple이나 bending constraint target으로 변환하는 확장이 가능하다.

\subsection*{Output}

\begin{equation}
f_i^{\mathrm{ext}}(t).
\end{equation}

출력은 simulation anchor별 external force field이다.
이제 force는 load sample domain이 아니라 solver vertex domain에 놓인다.

\subsection*{Effect}

force sampling result를 XPBD solver가 사용할 vertex force로 변환한다. 이 단계는 fixed load samples와 fixed simulation anchors를 연결한다.
여기서도 topology나 binding은 바뀌지 않는다.
Frame마다 바뀌는 것은 $f_l(t)$와 $f_i^{\mathrm{ext}}(t)$의 값뿐이다.

\section{Module 12: Proxy Dynamics Solver}

\subsection*{Input}

\begin{equation}
x_i(t),\quad v_i(t),\quad f_i^{\mathrm{ext}}(t),\quad m_i,\quad \theta_i,\quad C.
\end{equation}

입력은 현재 proxy vertex state와 외력, mass/material parameter, constraint set이다.
$C$는 stretch, bending, pin, attachment 같은 position-based constraints의 집합이다.
이 모듈은 Gaussian을 직접 움직이지 않고, proxy mesh의 physically plausible deformation을 먼저 계산한다.

\subsection*{Operation}

먼저 external force를 이용해 predicted velocity와 position을 계산한다.

\begin{align}
v_i^*&=v_i(t)+\Delta t\frac{f_i^{\mathrm{ext}}(t)}{m_i},\\
x_i^*&=x_i(t)+\Delta t\,v_i^*.
\end{align}

$v_i^*$와 $x_i^*$는 constraint projection 전의 tentative state이다.
외력은 이 단계에서 motion을 예측하고, 이후 XPBD constraints가 stretch, bending, pin 조건을 만족하도록 위치를 보정한다.

stretch constraint:

\begin{equation}
C_{ik}^{\mathrm{stretch}}(x)=\|x_i-x_k\|-\ell_{ik}^0,
\qquad \ell_{ik}^0=\|x_i^0-x_k^0\|.
\end{equation}

Stretch constraint는 edge length가 rest length $\ell_{ik}^0$에서 너무 벗어나지 않게 한다.
Cloth나 paper-like object에서 면이 비현실적으로 늘어나는 것을 막는 기본 constraint이다.

bending constraint:

\begin{equation}
C_e^{\mathrm{bend}}(x)=\theta_e(x)-\theta_e^0.
\end{equation}

Bending constraint는 adjacent faces 사이의 dihedral angle을 rest angle $\theta_e^0$ 근처로 유지한다.
이 항이 약하면 천처럼 잘 휘고, 강하면 종이비행기나 연의 rib처럼 형태를 유지한다.

pin constraint:

\begin{equation}
C_i^{\mathrm{pin}}(x)=x_i-x_i^{\mathrm{target}}.
\end{equation}

Pin constraint는 support, handle, attachment point를 target position에 묶는다.
Hard pin과 soft pin의 차이는 이 constraint의 compliance를 0에 가깝게 둘지, 유한한 값으로 둘지에 있다.

XPBD update for a generic constraint:

\begin{align}
\Delta\lambda&=-\frac{C(x)+\widetilde{\alpha}\lambda}{\sum_iw_i\|\nabla_{x_i}C\|^2+\widetilde{\alpha}},\\
x_i&\leftarrow x_i+w_i\nabla_{x_i}C\Delta\lambda,
\end{align}

$\lambda$는 constraint multiplier이고, $\Delta\lambda$는 이번 projection iteration에서 추가되는 correction이다.
$\nabla_{x_i}C$는 constraint를 줄이기 위해 vertex $i$를 어느 방향으로 움직여야 하는지 알려준다.
XPBD는 compliance $\alpha$를 통해 stiffness를 time step과 iteration 수에 덜 민감하게 다룰 수 있어 interactive simulation에 적합하다.

where:

\begin{equation}
w_i=\frac{1}{m_i},\qquad \widetilde{\alpha}=\frac{\alpha}{\Delta t^2}.
\end{equation}

$w_i$는 inverse mass이므로 무거운 vertex는 덜 움직이고, 가벼운 vertex는 더 많이 correction된다.
$\widetilde{\alpha}$는 time-step-scaled compliance이며, $\alpha=0$이면 거의 hard constraint처럼 동작한다.

velocity update:

\begin{equation}
v_i(t+1)=\frac{x_i(t+1)-x_i(t)}{\Delta t}.
\end{equation}

Constraint projection이 끝난 뒤 위치 차분으로 velocity를 갱신한다.
필요하면 velocity damping이나 constraint-aware velocity correction을 추가할 수 있지만, MVP에서는 이 기본 update로 시작한다.

\subsection*{Output}

\begin{equation}
M(t+1)=(V(t+1),E,F).
\end{equation}

출력은 다음 frame의 deformed proxy mesh이다.
Topology $E,F$는 그대로이고 vertex positions $V(t+1)$만 바뀐다.

\subsection*{Effect}

connected surface deformation, boundary constraints, material stiffness, damping, and wind-induced force propagation을 처리한다. Gaussian을 직접 particle로 움직이는 것보다 topology-aware response를 만들 수 있다.
이 모듈 덕분에 wind response는 개별 Gaussian의 독립 이동이 아니라, mesh connectivity와 material constraints를 따른 deformation으로 표현된다.
다음 Module 13은 이 proxy deformation을 다시 Gaussian rendering representation으로 전달한다.

\section{Module 13: Gaussian Attribute Transport}

\subsection*{Input}

\begin{equation}
M(t+1),\qquad \mathcal{B}_j=(T_j,b_j,R_{T_j}^0,\delta_j),\qquad g_j^0.
\end{equation}

입력은 새로 계산된 deformed proxy, Module 5에서 저장한 Gaussian binding, 그리고 canonical Gaussian state이다.
이 모듈은 simulation proxy와 rendering 3DGS를 다시 연결하는 runtime bridge이다.
Binding이 고정되어 있기 때문에 각 Gaussian은 자신이 붙은 triangle의 현재 frame을 통해 업데이트된다.

\subsection*{Operation}

triangle $T_j=(a,b,c)$의 current local frame $R_{T_j}(t+1)$을 계산한다. Relative rotation은:

\begin{equation}
\Delta R_{T_j}(t+1)=R_{T_j}(t+1)(R_{T_j}^{0})^\top.
\end{equation}

$R_{T_j}^0$는 rest triangle frame이고, $R_{T_j}(t+1)$는 deformed triangle frame이다.
$\Delta R_{T_j}$는 rest frame에서 current frame으로 가는 local rotation이다.
이 rotation이 있어야 Gaussian ellipsoid가 surface normal/tangent 변화에 맞춰 함께 회전한다.

Gaussian center update:

\begin{equation}
\mu_j(t+1)=b_{ja}x_a(t+1)+b_{jb}x_b(t+1)+b_{jc}x_c(t+1)+R_{T_j}(t+1)\delta_j.
\end{equation}

앞의 세 항은 barycentric coordinate로 triangle 위 material point를 따라가는 위치이다.
$R_{T_j}(t+1)\delta_j$는 optional local offset을 current frame 방향으로 회전시켜 더한다.
따라서 Gaussian이 surface 위 또는 근처의 같은 material 위치를 따라 움직인다.

covariance transport:

\begin{equation}
\Sigma_j(t+1)=\Delta R_{T_j}(t+1)\Sigma_j^0\Delta R_{T_j}(t+1)^\top.
\end{equation}

이 식은 Gaussian center만 translation하는 것이 아니라 anisotropic covariance frame도 함께 transport한다.
즉 leaf나 cloth surface가 회전하면 Gaussian ellipsoid의 얇은 normal 방향도 함께 회전한다.
이 처리가 빠지면 위치는 따라가더라도 splat orientation이 canonical 상태에 남아 surface와 어긋나는 artifact가 생길 수 있다.

선택적으로 face area stretch를 이용한 scale correction을 사용할 수 있다.

\begin{equation}
a_T(t)=\frac{\operatorname{area}(T(t))}{\operatorname{area}(T^0)}.
\end{equation}

$a_T(t)$는 triangle area가 rest 대비 얼마나 늘거나 줄었는지를 나타낸다.
이를 이용하면 surface stretch에 따라 Gaussian scale을 약간 키우거나 줄일 수 있다.
다만 scale을 직접 바꾸면 splat size popping이나 density 변화가 생길 수 있으므로, 초기 구현에서는 보수적으로 다룬다.

단, 초기 구현에서는 splat popping을 피하기 위해 center와 covariance rotation만 먼저 검증하는 것이 안전하다.

\subsection*{Output}

\begin{equation}
\mu_j(t+1),\quad \Sigma_j(t+1),\quad R_j(t+1),\quad s_j(t+1).
\end{equation}

Renderer가 covariance matrix를 직접 받지 않고 rotation/scale parameter를 요구하는 경우, $\Sigma_j(t+1)$를 다시 eigen-decomposition하거나 transported frame과 scale로 분해해 $R_j(t+1),s_j(t+1)$를 만든다.
Opacity와 appearance는 기본적으로 canonical 값을 유지하고, 필요한 correction은 다음 Module 14에서 처리한다.

\subsection*{Effect}

mesh deformation을 다시 3DGS rendering representation으로 변환한다. 이는 position-only update가 아니라 mean, anisotropic covariance, local frame을 포함하는 Gaussian attribute transport이다.
이 모듈이 논문의 persistent representation bridge에서 가장 직접적인 runtime 단계이다.
사용자 관점에서는 proxy mesh를 보지 않지만, Gaussian rendering 결과는 이 transport를 통해 wind-driven deformation을 반영한다.

\section{Module 14: SH / Appearance Correction}

\subsection*{Input}

\begin{equation}
c_j^0,\quad \Delta T_j(t),\quad \nabla x_{T_j}(t),\quad e_{T_j}(t),\quad W_{T_j}(t),\quad \theta_{T_j}.
\end{equation}

$c_j^0$는 canonical SH 또는 appearance coefficients이다.
$\Delta T_j(t)$는 Gaussian이 붙은 triangle/local frame의 변화, $\nabla x_{T_j}(t)$는 local deformation gradient나 strain-like feature이다.
$e_{T_j}(t)$와 $W_{T_j}(t)$는 exposure와 wind feature이고, $\theta_{T_j}$는 material feature이다.
이 모듈은 geometry transport만으로 설명되지 않는 view-dependent appearance 변화를 보정한다.

\subsection*{Operation}

canonical SH를 local frame 변화에 맞게 analytic correction하고, residual만 예측한다.

\begin{equation}
c_j(t)=R_{\mathrm{SH}}(\Delta T_j(t))c_j^0+\Delta c_j(t).
\end{equation}

$R_{\mathrm{SH}}(\Delta T_j(t))$는 local frame rotation이 SH coefficients에 미치는 analytic rotation이다.
Geometry가 회전하면 view-dependent lobes도 같은 방향으로 회전해야 하므로, 이 부분은 가능하면 closed-form correction으로 처리한다.
$\Delta c_j(t)$는 analytic rotation만으로 설명되지 않는 shading, thickness, exposure, deformation-induced appearance residual이다.

residual predictor:

\begin{equation}
\Delta c_j(t)=Q_\psi(c_j^0,\Delta T_j(t),\nabla x_{T_j}(t),e_{T_j}(t),W_{T_j}(t),\theta_{T_j}).
\end{equation}

$Q_\psi$는 작은 MLP나 lightweight module로 둘 수 있다.
입력 feature가 local binding/frame 기준으로 정의되면, asset 전체 좌표계가 바뀌어도 residual이 더 안정적으로 동작한다.
중요한 점은 SH residual이 geometry/dynamics를 대체하지 않고, runtime rendering fidelity를 보완한다는 것이다.

\subsection*{Output}

\begin{equation}
c_j(t).
\end{equation}

출력은 현재 frame에서 renderer에 넣을 appearance coefficients이다.
MVP에서는 $\Delta c_j(t)=0$으로 두고 analytic SH rotation만 검증하거나, degree-0 color는 고정하고 higher-order SH만 조정하는 단계적 구현도 가능하다.

\subsection*{Effect}

이 모듈은 deformation-aware appearance fidelity를 보완한다. 하지만 headline contribution이 아니라, geometry/dynamics pipeline을 완성하기 위한 supporting module로 두는 것이 좋다.
특히 runtime 구동을 위해 SH residual을 남겨두면, geometry transport 후에도 남는 view-dependent artifact를 가볍게 흡수할 수 있다.
다만 이 모듈의 성능을 주장하려면 geometry-only, analytic SH rotation only, residual 추가의 ablation이 필요하다.

\section{Module 15: 3DGS Rendering and Response Output}

\subsection*{Input}

\begin{equation}
G(t)=\{\mu_j(t),\Sigma_j(t),\alpha_j(t),c_j(t)\},\qquad \Pi_t.
\end{equation}

$G(t)$는 모든 runtime update가 반영된 dynamic Gaussian set이다.
각 Gaussian은 transported mean/covariance, opacity, corrected appearance를 갖는다.
$\Pi_t$는 현재 frame의 camera intrinsics/extrinsics이며, preview에서는 fixed camera나 turntable path, evaluation에서는 ground-truth camera path가 될 수 있다.

\subsection*{Operation}

standard 3DGS renderer를 사용한다.

\begin{equation}
I_t=\mathcal{R}_{\mathrm{3DGS}}(G(t),\Pi_t).
\end{equation}

이 단계에서는 새로운 renderer를 발명하는 것이 아니라, updated Gaussian state를 기존 3DGS renderer에 넣는다.
Renderer backend는 gsplat, diff-gaussian-rasterization, or CPU/debug renderer일 수 있지만 논리적으로는 같은 함수 $\mathcal{R}_{\mathrm{3DGS}}$로 본다.
따라서 contribution은 rendering algorithm 자체보다, renderer에 들어가기 전 dynamic $G(t)$를 어떻게 안정적으로 만드는지에 있다.

\subsection*{Output}

\begin{equation}
\boxed{\mathcal{I}=\{I_t\}_{t=1}^{T}}\qquad \text{or}\qquad \boxed{I_t}.
\end{equation}

단일 frame rendering이면 $I_t$가 출력이고, animation이나 simulation preview이면 frame sequence $\mathcal{I}$가 출력이다.
이 결과가 사용자가 직접 보는 최종 response이다.

export 가능한 dynamic representation은:

\begin{equation}
\boxed{\{G(t)\}_{t=1}^{T}}.
\end{equation}

$\{G(t)\}_{t=1}^{T}$를 저장하면 rendering 결과뿐 아니라 dynamic Gaussian asset 자체를 export할 수 있다.
이 representation은 다른 camera path로 다시 render하거나, downstream editing/evaluation에 사용할 수 있다.

\subsection*{Effect}

사용자는 최종적으로 high-quality 3DGS novel-view rendering response를 얻는다. 내부적으로는 mesh proxy가 simulation을 담당하지만, 입력과 출력은 여전히 3DGS asset 중심으로 유지된다.
이 점이 pipeline의 중요한 framing이다.
Mesh proxy는 visible product가 아니라 deformation scaffold이고, 최종 user-facing representation은 trained 3DGS와 그 dynamic response이다.

\section{Runtime Loop}

Asset setup과 training이 완료되어 다음이 준비되어 있다고 가정한다.

\begin{equation}
G^0,\quad M^0,\quad \mathcal{B},\quad Q_{\mathrm{load}},\quad P_\theta,\quad Q_\psi.
\end{equation}

$G^0$는 canonical Gaussian asset이고, $M^0$는 fixed-topology simulation proxy이다.
$\mathcal{B}$는 Gaussian-to-proxy binding table, $Q_{\mathrm{load}}$는 fixed material-space load samples이다.
$P_\theta$와 $Q_\psi$는 각각 force predictor와 optional SH residual predictor이며, analytic baseline만 사용할 경우에는 learned predictor 없이도 loop를 실행할 수 있다.
이 목록은 runtime loop가 시작되기 전에 asset setup 또는 training/preprocessing 단계에서 준비되어 있어야 하는 persistent state이다.

각 frame에서의 runtime loop는 다음과 같다.

\begin{enumerate}[leftmargin=1.5em]
    \item Load sample 위치 계산:
    \[
    x_l(t)=\sum_{i\in T_l}b_{li}x_i(t).
    \]
    \item Wind 계산:
    \[
    W_l(t)=W(x_l(t),t).
    \]
    \item Exposure 계산:
    \[
    e_l(t)=\operatorname{sample}(S_t,\pi_w(x_l(t))).
    \]
    \item Effective wind 계산:
    \[
    u_l^{\mathrm{eff}}(t)=e_l(t)W_l(t)-v_l(t).
    \]
    \item Force 예측:
    \[
    (f_l(t),\tau_l(t))=P_\theta(z_l(t),u_l^{\mathrm{eff}}(t),e_l(t),\theta_l).
    \]
    \item Vertex force 누적:
    \[
    f_i^{\mathrm{ext}}(t)=\sum_{l:i\in T_l}b_{li}f_l(t).
    \]
    \item XPBD/PBD solver:
    \[
    x(t+1)=\Phi_{\mathrm{XPBD}}(x(t),v(t),f^{\mathrm{ext}}(t),C,\theta).
    \]
    \item Gaussian geometry update:
    \[
    \mu_j(t+1)=\sum_{k\in T_j}b_{jk}x_k(t+1)+R_{T_j}(t+1)\delta_j,
    \]
    \[
    \Sigma_j(t+1)=\Delta R_{T_j}(t+1)\Sigma_j^0\Delta R_{T_j}(t+1)^\top.
    \]
    \item Gaussian appearance update:
    \[
    c_j(t+1)=R_{\mathrm{SH}}(\Delta T_j(t+1))c_j^0+Q_\psi(\cdots).
    \]
    \item Render:
    \[
	    I_{t+1}=\mathcal{R}_{\mathrm{3DGS}}(G(t+1),\Pi_{t+1}).
	    \]
\end{enumerate}

이 loop에서 매 frame 바뀌는 것은 wind query, exposure value, load force, proxy vertex positions, transported Gaussian attributes이다.
반대로 proxy topology, Gaussian binding, simulation anchor index, load sample material coordinate는 바뀌지 않는다.
이 separation이 temporal stability와 realtime implementation의 핵심이다.

\section{Minimum Development Version}

구현 복잡도를 낮추려면 다음 최소 버전부터 시작한다.

\begin{enumerate}[leftmargin=1.5em]
    \item trained 3DGS input 또는 synthetic Gaussians
    \item fixed-topology mesh proxy, 처음에는 known mesh 사용 가능
    \item Gaussian-to-triangle binding
    \item procedural wind load
    \item XPBD/PBD mesh deformation
    \item Gaussian mean/covariance transport
    \item canonical SH 또는 analytic SH correction
    \item 3DGS rendering output
\end{enumerate}

이 최소 버전은 논문 최종 시스템의 모든 장치를 한 번에 구현하려는 계획이 아니다.
먼저 ``mesh가 움직이면 Gaussian mean과 covariance가 안정적으로 따라가는가''를 확인하는 loop이다.
따라서 exposure estimator, dense Blender supervision, learned force predictor, SH residual은 이 단계에서 없어도 된다.

이 최소 loop는 다음과 같다.

\begin{equation}
G^0\rightarrow M^0\rightarrow \mathcal{B}\rightarrow M(t)\rightarrow G(t)\rightarrow I_t.
\end{equation}

여기서 $G^0\rightarrow M^0$는 proxy construction, $M^0\rightarrow\mathcal{B}$는 Gaussian binding, $M(t)\rightarrow G(t)$는 attribute transport를 의미한다.
초기 implementation에서는 known mesh와 synthetic Gaussians를 사용해도 좋다.
핵심 성공 기준은 visual plausibility보다 numeric/structural correctness이다: binding이 고정되는가, covariance가 triangle rotation을 따라가는가, solver topology가 유지되는가, renderer가 updated Gaussian을 정상적으로 받는가.

나중에 추가할 모듈은 Gaussian-opacity exposure, Blender dense force projection, learned force predictor, SH residual, material parameter calibration, projection-risk-based load sample refinement이다.
이 순서는 위험도가 낮은 core bridge부터 확인한 뒤, force realism과 appearance fidelity를 점진적으로 올리는 개발 전략이다.

\section{Module Hierarchy for Paper Framing}

현재 pipeline은 구성요소가 많기 때문에, 논문에서 모든 것을 같은 수준의 contribution으로 쓰면 안 된다. 권장 hierarchy는 다음이다.

\begin{center}
\begin{tabularx}{0.96\linewidth}{L{0.22\linewidth}Y}
\toprule
위계 & 모듈 \\
\midrule
Main contribution & fixed-topology mesh-proxy scaffold for wind-deformable 3DGS \\
Main contribution & persistent Gaussian-to-proxy binding and attribute transport \\
Secondary contribution & force-space wind deformation on proxy, with Blender dense-force supervision \\
Supporting module & Gaussian-opacity exposure estimator \\
Supporting module & conservation-preserving dense-to-proxy force projection \\
Supporting module & deformation-aware SH residual correction \\
Implementation detail & anchor/load sample scoring, material parameter regularization, projection-risk refinement \\
\bottomrule
\end{tabularx}
\end{center}

이 hierarchy의 의도는 시스템을 약하게 보이게 하려는 것이 아니라, reviewer가 핵심 claim을 놓치지 않게 하는 것이다.
Main contribution은 3DGS와 deformable proxy 사이의 persistent bridge이고, secondary/supporting modules는 그 bridge를 더 안정적이고 설득력 있게 만드는 장치이다.
특히 SH residual과 dense-force projection은 중요하지만, 그것들이 논문의 주장을 대신하지 않도록 서술해야 한다.

\section{Design Dependencies}

\begin{itemize}[leftmargin=1.4em]
    \item \textbf{Proxy mesh} provides topology, constraints, local frames, and force propagation.
    \item \textbf{Persistent binding} enables stable Gaussian attribute transport without per-frame rebinding.
    \item \textbf{Fixed load samples} allow time-varying wind exposure and force values without changing anchors.
    \item \textbf{Dense-to-proxy projection} compresses Blender dense supervision while preserving net force and wrench.
    \item \textbf{SH residual} improves rendering fidelity but should not dominate the paper narrative.
\end{itemize}

이 dependency를 거꾸로 읽으면 구현 순서도 된다.
Proxy mesh와 persistent binding이 먼저 안정되어야 wind force나 SH residual을 얹었을 때 문제가 어디에서 생겼는지 구분할 수 있다.
Fixed load samples는 wind exposure가 시간에 따라 바뀌어도 correspondence를 유지하게 해주고, dense-to-proxy projection은 학습/평가 단계에서 sparse load representation을 정당화한다.

\section{Suggested Coding-Agent Prompt}

\begin{lstlisting}
Use the pipeline note `3dgs_wind_response_pipeline_2026-05-12.tex` as the authoritative module-order reference.
Revise the idea sketch so that the method is presented as a Persistent Mesh-Proxy Wind Response Pipeline for 3DGS.
Separate asset setup, training/preprocessing, and runtime modules.
Make clear that runtime does not perform remeshing, Gaussian rebinding, or load-sample relocation.
Use simulation mesh vertices as fixed simulation anchors and fixed material-space quadrature points as load samples.
Present Gaussian-to-proxy binding as the key bridge for transporting mean, covariance, local frame, and appearance.
Keep Blender dense-force supervision and conservation-preserving force projection as supporting/training modules.
Keep SH residual correction as a supporting appearance module, not the headline contribution.
\end{lstlisting}

\section{Decision Log Entry}

\subsection*{2026-05-12 Pipeline Decomposition}

\textbf{Decision.} Record the method as the \textbf{Persistent Mesh-Proxy Wind Response Pipeline for 3DGS}. The pipeline starts from a trained 3DGS asset, constructs a fixed-topology simulation proxy, sets fixed simulation anchors and fixed material-space load samples, binds Gaussians persistently to proxy triangles, computes wind loads at runtime, solves proxy dynamics, transports Gaussian attributes, and renders dynamic 3DGS frames.

\textbf{Reason.} The current idea sketch contains many modules. Decomposing them by execution phase and dependency clarifies that only a compact subset runs every frame. This also keeps the paper framing focused on the representation bridge rather than presenting every supporting mechanism as a separate contribution.

\textbf{Next.} Rewrite the method section around four high-level blocks: persistent mesh-proxy scaffold, Gaussian-to-proxy binding and attribute transport, force-space proxy wind dynamics, and rendering/appearance correction. Move detailed anchor/load sample scoring and projection-risk refinement to implementation details or supplementary material.

\section{One-Paragraph Method Summary}

Given a trained 3DGS asset $G^0$, we first construct a fixed-topology simulation proxy $M^0$ and bind each Gaussian to a proxy triangle using persistent barycentric coordinates and a local frame. At runtime, a wind field and Gaussian/proxy-based exposure estimator produce aerodynamic loads on fixed material-space load samples. These loads are projected to simulation anchors and integrated by an XPBD solver to obtain the deformed proxy $M(t)$. The proxy deformation is then transported back to the Gaussian means, covariances, and local frames through the stored bindings. Finally, canonical SH coefficients are analytically corrected and optionally refined by a lightweight residual network before standard 3DGS rendering produces the output frame.

\bibliographystyle{plain}
\bibliography{refs}

\end{document}
